\section{Použité vývojové prostriedky}
\label{chapter:vyvojove-prostriedky}

V~tejto kapitole popisujeme technológie použité pri
vývoji aplikácie a~zdôvodňujeme ich voľbu. Aplikácia
pozostáva z~troch komunikujúcich častí -- časti
bežiacej na hodinkách, prepojovacej služby na
mobilnom telefóne bežiacej v~aplikácii Zepp
a~Android aplikácie zobrazujúcej stav zápasu. Okrem
samotnej aplikácie sme počas vývoja využívali aj
pomocné Python skripty na offline analýzu nameraných
údajov. Detaily o~tom, ako sú tieto technológie
nasadené v~konkrétnych komponentoch nášho riešenia,
uvádzame v~kapitole~\ref{chapter:implementacia}.

\subsection{Vývoj časti pre inteligentné hodinky}
\label{subsec:hodinky-tech}

\subsubsection{JavaScript a~Zepp OS Mini Program}

Aplikácie pre Zepp OS \cite{zepp_architecture} sa vyvíjajú výhradne
v~JavaScripte vo forme tzv. Mini Programov --
odľahčených aplikácií navrhnutých pre obmedzené
prostriedky inteligentných hodiniek. Mini Program
sa skladá z~troch logických celkov: \textbf{Device App}
bežiacej priamo na hodinkách, \textbf{Side Service}
bežiacej na mobilnom telefóne ako prepojovacia služba,
a~\textbf{Settings App} s~používateľským rozhraním
pre konfiguráciu. Voľba JavaScriptu nie je voliteľná;
ide o~jediný oficiálne podporovaný spôsob vývoja
pre túto platformu.

\subsubsection{Zepp OS Sensor API}

Prístup k~senzorom hodiniek poskytuje v~Zepp OS
modul \texttt{@zos/sensor}, ktorý sprístupňuje
akcelerometer a~gyroskop od verzie API~3.0
\cite{zepp_sensor_acc, zepp_sensor_gyro}.
Akcelerometer meria zrýchlenie zariadenia pozdĺž
troch ortogonálnych osí (X, Y, Z), pričom osi
X~a~Y sú rovnobežné s~displejom a~os~Z je naň
kolmá. Gyroskop meria uhlovú rýchlosť rotácie
taktiež v~troch osiach. Oba senzory fungujú
na princípe registrácie callback funkcie, ktorá
sa zavolá pri každej novej vzorke dát:

\begin{lstlisting}[language=JavaScript, caption=Príklad inicializácie akcelerometra v~Zepp OS]
import { Accelerometer, FREQ_MODE_NORMAL } from '@zos/sensor'

const accelerometer = new Accelerometer()
accelerometer.onChange(() => {
    console.log(accelerometer.getCurrent())
})
accelerometer.setFreqMode(FREQ_MODE_NORMAL)
accelerometer.start()
\end{lstlisting}

\subsubsection{Node.js a~Zeus CLI}

Node.js~\cite{nodejs-definition} je prostredie na beh JavaScriptu mimo
webového prehliadača, postavené na engine~V8
od spoločnosti Google.
V~našom projekte slúži ako prerekvizita pre
vývojový nástroj Zeus~\Gls{cli}, distribuovaný
ako npm balík \texttt{@zeppos/zeus-cli} a~vyžadujúci
Node.js vo verzii 14.0.0 alebo vyššej. Zeus \Gls{cli}
zahŕňa celý vývojový cyklus aplikácie pre Zepp OS --
od vytvorenia nového projektu (\texttt{zeus create}),
cez lokálne ladenie v~simulátore (\texttt{zeus dev})
a~na reálnom zariadení (\texttt{zeus preview}), až
po zostavenie finálneho inštalačného balíčka
(\texttt{zeus build}) \cite{zepp_cli}. Pre správu
samotného balíka slúži štandardný správca balíčkov
\Gls{npm}, ktorý je súčasťou inštalácie Node.js
\cite{npm-definition}.

\subsubsection{Komunikácia cez Bluetooth Low Energy}

Hodinky a~telefón medzi sebou komunikujú bezdrôtovo
prostredníctvom \Gls{ble} \cite{zepp_bluetooth}.
Keďže natívne \Gls{ble} \Gls{api} v~Zepp OS je
pomerne nízkoúrovňové, platforma poskytuje pomocnú
knižnicu \texttt{MessageBuilder}, ktorá abstrahuje
odosielanie a~príjem správ medzi Device App a~Side
Service do jednoduchého rozhrania.

\subsection{Vývoj časti pre mobilný telefón}
\label{subsec:android-tech}

\subsubsection{Kotlin}

Pre vývoj Android aplikácie sme zvolili jazyk
Kotlin \cite{kotlin_android}, ktorý je od roku 2019 preferovaným jazykom
Googlu pre vývoj Android aplikácií. Oproti Jave prináša Kotlin
výrazne stručnejšiu syntax, zabudovanú ochranu
pred \textit{null pointer} výnimkami a~podporu pre
korutiny (\textit{coroutines}), ktoré zjednodušujú
prácu s~asynchrónnym kódom \cite{kotlin_overview}.

\subsubsection{Android Studio}

Android Studio~\cite{android_studio_intro} je oficiálne \Gls{ide} pre vývoj 
Android aplikácií. Je postavené na platforme IntelliJ IDEA od
spoločnosti JetBrains a~bolo predstavené na
konferencii Google~I/O v~roku 2013. Vďaka vstavanému
emulátoru, profilovaciemu nástroju a~build systému
Gradle pokrýva celý vývojový cyklus Android
aplikácie.

\subsection{Spracovanie a~analýza dát}
\label{subsec:python-tech}

\subsubsection{Python}

Python~\cite{python-definition} je interpretovaný, vysokoúrovňový programovací
jazyk s~dynamickým typovaním, ktorý kladie dôraz
na čitateľnosť a~jednoduchosť syntaxe. Vďaka rozsiahlemu
ekosystému knižníc pre prácu s~dátami sa stal
štandardom pre analýzu nameraných údajov vo
vedeckom prostredí. V~našom projekte slúži na
pomocné skripty pre kalibráciu intervalov hodnôt
jednotlivých gest -- detaily o~konkrétnom
použití uvádzame v~podsekcii
\ref{sec:python-skript}.